\documentclass[conference]{IEEEtran}
\usepackage[utf8]{inputenc}
\usepackage[T1]{fontenc}
\usepackage{amsmath,amssymb}
\usepackage{graphicx}
\usepackage{hyperref}
\usepackage{xcolor}
\usepackage{booktabs}

\begin{document}

\title{Trabalho 2 -- CIC0201: Cifragem RSA-OAEP e Assinatura Digital RSA-PSS com SHA3-256}

\author{
\IEEEauthorblockN{Rafael de Lima Pereira}
\IEEEauthorblockA{Matrícula: 242043277}
\and
\IEEEauthorblockN{Pedro Freitas França}
\IEEEauthorblockA{Matrícula: 232006028}
\and
\IEEEauthorblockN{Rian Kallebe da Silva Lisboa}
\IEEEauthorblockA{Matrícula: 242012000}
\and
\IEEEauthorblockN{Arthur Martins Pereira de Souza}
\IEEEauthorblockA{Matrícula: 241004499}
}

\maketitle

\begin{abstract}
Este trabalho apresenta uma implementação didática, em Python puro, do algoritmo RSA e dos
esquemas RSA-OAEP e RSA-PSS com SHA3-256, cobrindo geração de chaves, cifragem, assinatura
digital de arquivos e verificação de integridade. Toda a aritmética modular, o teste de
primalidade de Miller-Rabin e as codificações de padding foram implementados manualmente,
sem bibliotecas criptográficas externas na via principal. Os resultados experimentais,
obtidos por uma suíte com 46 testes automatizados e por bibliotecas independentes usadas
apenas para verificação cruzada, confirmam a corretude do esquema e sua capacidade de
rejeitar adulterações no conteúdo, na assinatura e na chave.
\end{abstract}

\begin{IEEEkeywords}
RSA, OAEP, PSS, SHA3-256, assinatura digital, segurança computacional
\end{IEEEkeywords}

\section{Introdução}

A disciplina de Segurança Computacional (CIC0201), inserida no contexto mais amplo de redes
de computadores, trata da garantia de propriedades como confidencialidade, integridade e
autenticidade em comunicações e no armazenamento de dados. Sistemas de criptografia de chave
pública, como o RSA, são fundamento de protocolos amplamente usados em redes, a exemplo do
TLS e da assinatura de pacotes de software, permitindo que partes que nunca compartilharam um
segredo estabeleçam comunicação segura ou verifiquem a origem de um arquivo.

Este trabalho tem como objetivo consolidar, por meio de implementação manual, o entendimento
dos mecanismos internos do RSA e dos esquemas de preenchimento (\textit{padding}) que o tornam
seguro na prática: o OAEP (\textit{Optimal Asymmetric Encryption Padding}), usado para
cifragem, e o PSS (\textit{Probabilistic Signature Scheme}), usado para assinatura digital.
Diferentemente do uso de uma biblioteca pronta, a implementação dos algoritmos matemáticos e
das codificações a partir da especificação evidencia por que RSA sem preenchimento é inseguro
e como cada etapa dos esquemas OAEP e PSS contribui para as garantias de segurança
correspondentes.

O relatório está estruturado da seguinte forma: a Seção II apresenta a fundamentação teórica
sobre aritmética modular, geração de chaves RSA, OAEP e PSS; a Seção III descreve o ambiente
experimental e analisa os resultados obtidos com a suíte de testes e com a demonstração de
ponta a ponta; a Seção IV apresenta as conclusões técnicas; e a Seção V lista as referências
bibliográficas, que também indica o repositório público onde o código-fonte completo do
projeto está disponível~\cite{repo}.

\section{Fundamentação Teórica}

\subsection{Aritmética modular e geração de chaves RSA}

O RSA baseia-se na dificuldade computacional de fatorar o produto de dois números primos
grandes~\cite{menezes1996handbook}. A geração de chaves consiste em escolher dois primos $p$
e $q$ distintos, calcular o módulo $n = p \cdot q$ e a função totiente de Euler
$\varphi(n) = (p-1)(q-1)$, selecionar um expoente público $e$ coprimo com $\varphi(n)$ e
calcular o expoente privado $d$ como o inverso multiplicativo de $e$ módulo $\varphi(n)$, isto
é, $e \cdot d \equiv 1 \pmod{\varphi(n)}$.

Neste projeto, a exponenciação modular rápida foi implementada pelo método
\textit{square-and-multiply}, que calcula $\text{base}^{\text{exp}} \bmod n$ em tempo
$O(\log(\text{exp}))$ multiplicações, e o inverso modular foi obtido pelo Algoritmo Estendido
de Euclides. A geração dos primos $p$ e $q$ usa candidatos aleatórios de tamanho fixo (obtidos
com a fonte de aleatoriedade criptograficamente segura do sistema operacional, via
\texttt{secrets.randbits}) filtrados pelo teste probabilístico de primalidade de Miller-Rabin.
Esse teste explora a fatoração $n - 1 = 2^{s}d$, com $d$ ímpar, e testemunhas aleatórias $a$: se
nenhuma sequência de quadrados sucessivos de $a^{d}$ atinge $n-1$ antes de $a^{2^{s-1}d}$, o
número é declarado composto; caso contrário, é provavelmente primo. Com $r$ rodadas
independentes, a probabilidade de um composto ser aceito erroneamente é, no pior caso, no
máximo $4^{-r}$~\cite{menezes1996handbook}; o projeto usa 20 rodadas por padrão. Também são
calculados os parâmetros do Teorema Chinês do Resto ($d_p$, $d_q$, $q^{-1} \bmod p$), que
aceleram operações com a chave privada.

As chaves geradas (módulo mínimo de 2048 bits) são serializadas em um formato JSON próprio,
com os inteiros grandes representados em Base64 sobre uma codificação \textit{big-endian}. A
importação de chaves privadas revalida estruturalmente todos os parâmetros e as relações
algébricas ($p \cdot q = n$, $e \cdot d \equiv 1 \bmod \varphi(n)$, consistência de $d_p$,
$d_q$ e $q^{-1}$) e reaplica Miller-Rabin aos fatores, de modo a detectar arquivos de chave
corrompidos ou adulterados antes de qualquer operação criptográfica.

\subsection{RSA-OAEP: cifragem com preenchimento ótimo}

A aplicação direta do RSA a uma mensagem (RSA "de livro-texto") é determinística e
matematicamente estruturada, o que a torna vulnerável a ataques que exploram a maleabilidade
multiplicativa do esquema. O OAEP, especificado na RFC~8017~\cite{rfc8017}, resolve esse
problema combinando a mensagem com uma semente aleatória antes da exponenciação pública,
usando uma função geradora de máscara (MGF1) construída sobre uma função de hash -- neste
trabalho, SHA3-256~\cite{fips202}. A cifragem calcula o hash de um rótulo (\textit{label}),
monta um bloco de dados com preenchimento de zeros e um separador, mascara esse bloco com uma
semente aleatória via MGF1 e, em seguida, mascara a própria semente com o hash do bloco
mascarado, produzindo uma cifragem probabilística: mensagens iguais cifradas com a mesma chave
produzem criptogramas diferentes a cada execução. A decifragem inverte as máscaras, recalcula
o hash do rótulo e verifica cada zero de preenchimento e o byte separador com comparação de
tempo constante (\texttt{hmac.compare\_digest}) antes de aceitar a mensagem recuperada,
rejeitando com uma única mensagem de erro genérica qualquer falha estrutural.

\subsection{RSA-PSS: assinatura digital probabilística}

O PSS, também especificado na RFC~8017~\cite{rfc8017}, é um esquema de assinatura que combina
o hash da mensagem com um \textit{salt} aleatório antes de codificar um bloco de dados
mascarado por MGF1, finalizado com um byte de traço (\textit{trailer}) fixo (\texttt{0xbc}).
Como o \textit{salt} é gerado aleatoriamente a cada assinatura, duas assinaturas da mesma
mensagem com a mesma chave são diferentes entre si, mas ambas válidas. A verificação recalcula
o hash da mensagem, extrai o \textit{salt} do bloco de dados desmascarado, reconstrói o hash
esperado e o compara ao hash embutido na assinatura com \texttt{hmac.compare\_digest}. Diferente
do OAEP, cuja finalidade é confidencialidade, o PSS provê integridade e autenticidade
relativas à posse da chave privada correspondente, sem ocultar o conteúdo assinado.

\subsection{Comparação com Ed25519}

Como contraponto ao RSA-PSS, o esquema de assinatura Ed25519 (RFC~8032~\cite{rfc8032}), baseado
no problema do logaritmo discreto sobre uma curva de Edwards, produz assinaturas de 64 bytes e
chaves públicas de 32 bytes, ordens de grandeza menores que os 256 bytes de assinatura e chave
usados pelo perfil RSA de 2048 bits deste trabalho. Além disso, o Ed25519 deriva o
\textit{nonce} de assinatura deterministicamente a partir da chave e da mensagem, dispensando
uma nova fonte de aleatoriedade de alta qualidade a cada assinatura -- ainda que a geração da
chave em si continue exigindo aleatoriedade segura. O tamanho do módulo RSA não é diretamente
comparável ao tamanho de chave de curvas elípticas em termos de nível de segurança
equivalente~\cite{ferguson2010cryptography}.

\section{Ambiente Experimental e Análise de Resultados}

\subsection{Descrição do Cenário}

Os experimentos foram executados em uma máquina com processador AMD Ryzen 5 5500 (12 threads
lógicas), 14~GiB de memória RAM, sistema operacional Linux (kernel 7.0.0, distribuição baseada
em Ubuntu) e interpretador Python 3.14.4. A implementação utiliza exclusivamente a biblioteca
padrão do Python (\texttt{hashlib} para SHA3-256, \texttt{secrets} para aleatoriedade
criptograficamente segura, \texttt{base64} e \texttt{json} para serialização) na via principal
de geração de chaves, cifragem, assinatura e verificação, coerente com os conceitos descritos
na Seção~II: nenhuma rotina de aritmética modular, MGF1, OAEP ou PSS é delegada a terceiros. O
projeto está organizado em três módulos executáveis (\texttt{rsa\_keygen.py},
\texttt{rsa\_crypto.py} e \texttt{signed\_file.py}) e uma suíte de seis arquivos de teste sob o
diretório \texttt{tests/}, executada com o executor de testes padrão da linguagem
(\texttt{unittest}). Duas bibliotecas de terceiros (\texttt{cryptography} e
\texttt{pycryptodome}) são usadas apenas em testes de interoperabilidade opcionais, isolados em
um ambiente virtual Python, para checar de forma independente as saídas PSS e OAEP produzidas
pela implementação do grupo -- essas bibliotecas nunca participam da implementação principal.

A interface de linha de comando de \texttt{signed\_file.py} expõe três operações: \texttt{keygen}
(gera e exporta um par de chaves RSA de 2048 bits), \texttt{sign} (lê um arquivo, assina seus
bytes com RSA-PSS e grava uma estrutura JSON assinada) e \texttt{verify} (importa uma chave
pública e valida a estrutura assinada), retornando códigos de saída distintos para sucesso (0),
assinatura ou estrutura inválida (1) e falhas operacionais como chave ou arquivo inválido (2).

\subsection{Análise de Resultados}

\paragraph{Geração de chaves e round-trip}
A geração de um par de chaves RSA de 2048 bits, incluindo a busca por dois primos de 1024 bits
aprovados em 20 rodadas de Miller-Rabin, completou em cerca de 0,39\,s neste ambiente. A
exportação para JSON, seguida da reimportação e da checagem das relações RSA/CRT, preservou
integralmente os parâmetros originais, confirmando o \textit{round-trip} de serialização.
Um teste de corrupção proposital (remoção do campo \texttt{e} da chave pública exportada) foi
corretamente rejeitado pelo importador antes de qualquer operação criptográfica.

\paragraph{Demonstração de assinatura e verificação}
Na demonstração de ponta a ponta, a assinatura RSA-PSS do arquivo \texttt{README.md}
(6014 bytes) levou aproximadamente 0,21\,s, produzindo uma estrutura JSON assinada de
8539 bytes; a verificação subsequente, dominada pela exponenciação modular com o expoente
público pequeno, levou apenas 0,04\,s e reportou
\textit{"Arquivo íntegro: assinatura RSA-PSS válida para a chave pública fornecida"}. Ao
substituir a chave pública por uma gerada de forma independente, a mesma estrutura assinada
foi corretamente rejeitada com a mensagem \textit{"Assinatura inválida: arquivo, estrutura ou
chave não correspondem"} e código de saída 1, confirmando que a verificação está de fato
vinculada à chave pública fornecida e não apenas ao formato da estrutura.

\paragraph{Suíte automatizada de testes}
A suíte completa, com 46 casos de teste distribuídos em seis arquivos, foi executada em
5,3\,s (um teste foi ignorado por depender de uma biblioteca opcional não instalada nesta
execução). Os testes de \texttt{test\_keygen.py} comparam a aritmética modular e o Euclides
estendido contra os operadores nativos do Python, validam parâmetros inválidos antes da
geração aleatória e verificam o teste de Miller-Rabin com testemunhas controladas, incluindo um
pseudoprimo forte conhecido. Os testes de \texttt{test\_oaep.py} cobrem cifragem/decifragem no
limite de tamanho de mensagem, rótulo incorreto, criptogramas adulterados ou fora do intervalo
$[0, n)$, padding malformado e um módulo de 2049 bits (não alinhado a um número inteiro de
bytes). Os testes de \texttt{test\_verification.py} alteram isoladamente um byte do conteúdo,
da assinatura e do módulo público, trocam a chave pública por outra válida, testam Base64 não
canônico e campos JSON duplicados, e confirmam que o \textit{salt} aleatório do PSS produz
assinaturas diferentes a cada execução sobre a mesma mensagem -- todos os casos de adulteração
foram corretamente rejeitados, e um caso válido foi mantido na suíte para evitar que um
verificador que rejeitasse tudo indevidamente passasse nos testes. Os testes de
\texttt{test\_private\_key.py} cobrem a importação de chaves privadas com JSON malformado,
campos ausentes, fatores compostos com equações algébricas consistentes (o que exercita
especificamente a repetição do teste de primalidade sobre os fatores) e alteração de cada
parâmetro numérico isoladamente. Por fim, os testes de interoperabilidade em
\texttt{test\_interop.py} confirmam que a biblioteca \texttt{cryptography} aceita como válida
uma assinatura PSS produzida pela implementação manual do grupo, dando evidência adicional de
conformidade com a especificação RFC~8017 além da suíte interna.

\paragraph{Ajustes e problemas encontrados}
Durante o desenvolvimento, a principal dificuldade tratada foi a validação estrita da
importação de chaves privadas: verificações puramente algébricas isoladas (por exemplo, apenas
$p \cdot q = n$) não garantem que $p$ e $q$ sejam de fato primos, permitindo que fatores
compostos artificialmente consistentes passassem despercebidos; a solução adotada foi reaplicar
Miller-Rabin aos fatores importados. Outro ajuste necessário foi tratar corretamente módulos
RSA cujo tamanho em bits não é múltiplo de 8 (como no teste de 2049 bits), o que exige atenção
ao número de bits usados no mascaramento do bloco de dados do PSS e do OAEP para não vazar
informação por bits fora do intervalo esperado. Também foi necessário padronizar as mensagens
de erro de decifragem OAEP para um único texto genérico, evitando que motivos de falha
distintos (rótulo incorreto, padding malformado, criptograma fora do intervalo) pudessem ser
diferenciados por um atacante observando apenas a saída da função.

\section{Conclusões}

A implementação manual de RSA, OAEP e PSS permitiu observar concretamente por que a
exponenciação RSA crua é insegura e como cada etapa dos esquemas de preenchimento contribui
para confidencialidade (OAEP) ou para integridade e autenticidade (PSS). A suíte de 46 testes
automatizados, incluindo casos de adulteração de conteúdo, assinatura e chave, bem como
verificação cruzada com bibliotecas independentes, deu evidência consistente de corretude do
esquema implementado dentro do perfil escolhido (SHA3-256, módulos de ao menos 2048 bits). Ao
mesmo tempo, a análise de segurança do projeto deixa claro que a implementação é didática: não
há garantias de tempo constante na exponenciação modular em Python, a chave privada é
armazenada sem cifragem, e a validação matemática da assinatura não substitui a autenticação
independente da chave pública nem prova a identidade de quem assina.

\begin{thebibliography}{9}

\bibitem{rfc8017}
K.~Moriarty, B.~Kaliski, J.~Jonsson, and A.~Rusch, ``PKCS \#1: RSA Cryptography Specifications
Version 2.2,'' RFC~8017, IETF, Nov. 2016. [Online]. Available: \url{https://www.rfc-editor.org/rfc/rfc8017}

\bibitem{fips202}
National Institute of Standards and Technology, ``SHA-3 Standard: Permutation-Based Hash and
Extendable-Output Functions,'' FIPS PUB 202, Aug. 2015. [Online]. Available:
\url{https://csrc.nist.gov/pubs/fips/202/final}

\bibitem{rfc8032}
S.~Josefsson and I.~Liusvaara, ``Edwards-Curve Digital Signature Algorithm (EdDSA),'' RFC~8032,
IETF, Jan. 2017. [Online]. Available: \url{https://www.rfc-editor.org/rfc/rfc8032}

\bibitem{menezes1996handbook}
A.~J. Menezes, P.~C. van Oorschot, and S.~A. Vanstone, \emph{Handbook of Applied Cryptography}.
Boca Raton, FL: CRC Press, 1996.

\bibitem{ferguson2010cryptography}
N.~Ferguson, B.~Schneier, and T.~Kohno, \emph{Cryptography Engineering: Design Principles and
Practical Applications}. Indianapolis, IN: Wiley, 2010.

\bibitem{fips186}
National Institute of Standards and Technology, ``Digital Signature Standard (DSS),''
FIPS PUB 186-5, Feb. 2023. [Online]. Available:
\url{https://csrc.nist.gov/pubs/fips/186/5/final}

\bibitem{repo}
R.~L. Pereira, P.~F. França, R.~K. da~S. Lisboa, and A.~M. P. de~Souza, ``Seguranca-Projeto2:
Sistema de assinatura digital e verificação de arquivos (RSA, OAEP e PSS com SHA3-256),''
repositório de código-fonte, GitHub, 2026. [Online]. Available:
\url{https://github.com/KallebeLisboa/Seguranca-Projeto2}

\end{thebibliography}

\end{document}
